% Source: physical pp. 364--366 (printed pp. 341--343), beginning at
% the section heading on p. 364 and ending immediately before section
% 31.6 on p. 366.
% Coverage: Eqs. (31.5.1)--(31.5.19), one unnumbered display, and no
% footnotes, citations, figures, tables, or typographic dividers.

\subsection{Local Supersymmetry Transformations}

As a last step before considering the effects of higher order in \(G\) on
supersymmetry transformation rules, we will now complete our discussion
of the transformation rules of the physical components of the graviton
superfield \(H_\mu\) and other superfields to lowest order in \(G\).

Let us first note the form that these transformation rules take when
expressed in terms of the physical fields \(h_{\mu\nu}\), \(\psi_\mu\),
\(b_\mu\), \(s\), and \(p\) in a `Wess--Zumino' gauge in which
\begin{equation}
  C_\mu^H=\omega_\mu^H=V_{\mu\nu}^H-V_{\nu\mu}^H=0\,.
  \tag{31.5.1}\label{eq:31.5.1}
\end{equation}
Using the identification \eqref{eq:31.1.25}--\eqref{eq:31.1.27},
\eqref{eq:31.1.29}, and \eqref{eq:31.1.11} of the physical fields with
components of the superfield \(H_\mu\), together with the general
transformation rules \eqref{eq:26.2.11}--\eqref{eq:26.2.17} and the
`gauge' conditions \eqref{eq:31.5.1}, we find the transformation rules
\begin{equation}
  \delta h_{\mu\nu}
  =\frac{1}{2}\bigl(\bar\alpha
    [\gamma_\mu\psi_\nu+\gamma_\nu\psi_\mu]\bigr),
  \tag{31.5.2}\label{eq:31.5.2}
\end{equation}
\begin{equation}
  \delta\psi_\mu
  =\left[
      \frac{1}{2}[\gamma^\nu,\gamma^\lambda]\partial_\lambda h_{\mu\nu}
      +\partial_\mu h^\lambda{}_\lambda
      +2\ii\gamma_5b_\mu
      -\frac{2}{3}\ii\gamma_\mu\gamma_\rho\gamma_5b^\rho
      +\frac{2}{3}\gamma_\mu s
      -\frac{2}{3}\ii\gamma_\mu\gamma_5p
    \right]\alpha,
  \tag{31.5.3}\label{eq:31.5.3}
\end{equation}
\begin{equation}
  \delta s=\frac{1}{4}\bigl(\bar\alpha\gamma^\lambda L_\lambda\bigr),
  \tag{31.5.4}\label{eq:31.5.4}
\end{equation}
\begin{equation}
  \delta p=-\frac{1}{4}i
    \bigl(\bar\alpha\gamma_5\gamma^\lambda L_\lambda\bigr),
  \tag{31.5.5}\label{eq:31.5.5}
\end{equation}
\begin{equation}
  \delta b_\mu
  =\frac{3}{4}\ii\bigl(\bar\alpha\gamma_5L_\mu\bigr)
   -\frac{1}{4}i
    \bigl(\bar\alpha\gamma_5\gamma_\mu\gamma^\nu L_\nu\bigr),
  \tag{31.5.6}\label{eq:31.5.6}
\end{equation}
where \(L_\mu\) is given by Eq. \eqref{eq:31.2.8}:
\[
  L^\mu\equiv
  \ii\epsilon^{\mu\nu\kappa\rho}\gamma_5\gamma_\nu
  \partial_\kappa\psi_\rho\,.
\]
This transformation shifts the components \(C_\mu^H\), \(\omega_\mu^H\),
and \(V_{\mu\nu}^H-V_{\nu\mu}^H\) away from zero, by the amounts
\begin{equation}
  \delta C_\mu^H=0,\qquad
  \delta\omega_\mu^H=V_{\mu\nu}^H\gamma^\nu\alpha,
  \tag{31.5.7}\label{eq:31.5.7}
\end{equation}
\begin{equation}
  \delta[V_{\nu\mu}^H-V_{\mu\nu}^H]
  =\bigl(\bar\alpha
    [\gamma_\mu\lambda_\nu^H-\gamma_\nu\lambda_\mu^H]\bigr).
  \tag{31.5.8}\label{eq:31.5.8}
\end{equation}
We can then return to the gauge satisfying Eq. \eqref{eq:31.5.1} by
making a suitable gauge transformation \(H_\mu\longrightarrow
H_\mu+\Delta_\mu\), where \(\Delta_\mu\) is a superfield of the form
\eqref{eq:31.1.17}, \eqref{eq:31.1.18}, with components
\begin{equation}
  C_\mu^\Delta=0,\qquad
  \omega_\mu^\Delta=-V_{\mu\nu}^H\gamma^\nu\alpha,
  \tag{31.5.9}\label{eq:31.5.9}
\end{equation}
\begin{equation}
  V_{\nu\mu}^\Delta-V_{\mu\nu}^\Delta
  =\bigl(\bar\alpha
    [\gamma_\nu\lambda_\mu^H-\gamma_\mu\lambda_\nu^H]\bigr).
  \tag{31.5.10}\label{eq:31.5.10}
\end{equation}

So far, this has been for global supersymmetry transformations, with
\(\alpha\) an infinitesimal constant Majorana spinor. At least in lowest
order, this symmetry can be easily extended to \emph{local}
supersymmetry transformations, with \(\alpha(x)\) given an arbitrary
dependence on \(x^\mu\). According to Eq. \eqref{eq:26.7.11}, under such
a transformation the matter action undergoes the change
\begin{equation}
  \delta\int \dd^4x\,\mathcal{L}_M
  =-\int \dd^4x\,
    \bigl(\bar S^\mu(x)\partial_\mu\alpha(x)\bigr).
  \tag{31.5.11}\label{eq:31.5.11}
\end{equation}
Inspection of Eq. \eqref{eq:31.2.17} shows that this change in the action
is cancelled if we add an inhomogeneous term
\((2/\kappa)\partial_\mu\alpha(x)\) to the right-hand side of
Eq. \eqref{eq:31.5.3}, so that the change in the gravitino field is now
\begin{multline}
  \delta\psi_\mu(x)
  =(2/\kappa)\partial_\mu\alpha(x)
  +\left[
      \frac{1}{2}[\gamma^\nu,\gamma^\lambda]
        \partial_\lambda h_{\mu\nu}(x)
      +\partial_\mu h^\lambda{}_\lambda(x)
      +2\ii\gamma_5b_\mu(x)\right.\\
  \left.
      -\frac{2}{3}\ii\gamma_\mu\gamma_\rho\gamma_5b^\rho(x)
      +\frac{2}{3}\gamma_\mu s(x)
      -\frac{2}{3}\ii\gamma_\mu\gamma_5p(x)
    \right]\alpha(x).
  \tag{31.5.12}\label{eq:31.5.12}
\end{multline}

It is useful to rewrite Eq. \eqref{eq:31.5.12} in a form that makes the
generalization to general coordinates more transparent. First, we note
that the term \(\partial_\mu h^\lambda{}_\lambda\alpha\) on the
right-hand side of Eq. \eqref{eq:31.5.12} may be eliminated by replacing
the parameter \(\alpha(x)\) with
\begin{equation}
  \widetilde\alpha
  \equiv(\operatorname{Det}g)^{1/4}\alpha
  \simeq\alpha+\frac{1}{2}\kappa h^\lambda{}_\lambda\alpha.
  \tag{31.5.13}\label{eq:31.5.13}
\end{equation}
Dropping the tilde, to zeroth order in \(\kappa\) Eq. \eqref{eq:31.5.12}
then becomes
\begin{multline}
  \delta\psi_\mu(x)
  =(2/\kappa)\partial_\mu\alpha(x)
  +\left[
      \frac{1}{2}[\gamma^\nu,\gamma^\lambda]
        \partial_\lambda h_{\mu\nu}(x)
      +2\ii\gamma_5b_\mu(x)\right.\\
  \left.
      -\frac{2}{3}\ii\gamma_\mu\gamma_\rho\gamma_5b^\rho(x)
      +\frac{2}{3}\gamma_\mu s(x)
      -\frac{2}{3}\ii\gamma_\mu\gamma_5p(x)
    \right]\alpha(x).
  \tag{31.5.14}\label{eq:31.5.14}
\end{multline}

We can express this in terms of the covariant derivative of \(\alpha(x)\),
which in general coordinates takes the form
\begin{equation}
  D_\mu\alpha(x)
  =\partial_\mu\alpha(x)
   +\frac{1}{2}\ii\mathscr{J}_{bc}\omega_\mu^{bc}(x)\alpha(x),
  \tag{31.5.15}\label{eq:31.5.15}
\end{equation}
where \(\mathscr{J}_{bc}\) is the matrix \eqref{eq:5.4.6} representing
the generator of Lorentz transformations in the Dirac representation
\begin{equation}
  \mathscr{J}^{bc}
  \equiv-\frac{\ii}{4}[\gamma^b,\gamma^c],
  \tag{31.5.16}\label{eq:31.5.16}
\end{equation}
and \(\omega_\mu^{bc}(x)\) is the spin connection,
\begin{equation}
  \omega_\mu^{bc}
  =e^b{}_\lambda e^c{}_{\nu;\mu}g^{\lambda\nu}
  =e^b{}_\lambda\frac{\partial e^c{}_\nu}{\partial x^\mu}g^{\lambda\nu}
   -\Gamma^\rho_{\nu\mu}e^b{}_\lambda e^c{}_\rho g^{\lambda\nu}.
  \tag{31.5.17}\label{eq:31.5.17}
\end{equation}
Using the weak field approximations \eqref{eq:31.1.5},
\eqref{eq:31.1.10}, and \eqref{eq:31.1.11}, together with the gauge
condition \(\phi_{\mu\nu}=\phi_{\nu\mu}\), and in this approximation
again ignoring the difference between local Lorentz indices \(a,b\),
etc.\ and spacetime indices \(\mu,\nu\), etc., this gives
\begin{equation}
  D_\mu\alpha(x)
  \simeq\partial_\mu\alpha(x)
  +\frac{1}{4}\kappa[\gamma^\nu,\gamma^\lambda]
    \partial_\lambda h_{\mu\nu}(x)\alpha(x).
  \tag{31.5.18}\label{eq:31.5.18}
\end{equation}
Thus the local supersymmetry transformation rule \eqref{eq:31.5.12} may
be written
\begin{multline}
  \delta\psi_\mu(x)
  =(2/\kappa)D_\mu\alpha(x)
  +\left[
      2\ii\gamma_5b_\mu(x)
      -\frac{2}{3}\ii\gamma_\mu\gamma_\rho\gamma_5b^\rho(x)\right.\\
  \left.
      +\frac{2}{3}\gamma_\mu s(x)
      -\frac{2}{3}\ii\gamma_\mu\gamma_5p(x)
    \right]\alpha(x).
  \tag{31.5.19}\label{eq:31.5.19}
\end{multline}
We see that in Wess--Zumino gauge the derivatives in the supersymmetry
transformation become covariant derivatives. In this sense, the
approach outlined here is similar to the de Wit--Freedman approach to
supersymmetric gauge theories, described in Section 27.8.

The transformation
\(\psi_\mu(x)\longrightarrow\psi_\mu(x)+(2/\kappa)\partial_\mu\alpha(x)\)
is a gauge transformation of the same type as Eq. \eqref{eq:31.1.13},
so it leaves the zeroth-order gravitino action
\(-\frac{1}{2}\int \dd^4x(\bar\psi_\mu L^\mu)\) invariant. The whole action
obtained from the Lagrangian density \eqref{eq:31.2.15} is therefore
invariant to order zero in \(\kappa\) under the local supersymmetry
transformations \eqref{eq:31.5.2}, \eqref{eq:31.5.4}--\eqref{eq:31.5.6},
\eqref{eq:31.5.12} (or \eqref{eq:31.5.19}) and matter superfield
transformations like \eqref{eq:26.7.15}. We conclude that
\emph{the combination of gravitation and supersymmetry automatically
leads to local supersymmetry}.
